\chapter{Trees and Boosting}\label{ch:ml:trees-and-boosting}

The firm's default gradient-boosted trees (300 trees of fifteen leaves) are fitted to ten years of a 300-stock panel
and score $-0.19\%$ out of sample. A random search of twenty configurations finds one that scores 0.39\% on the
validation months and 0.32\% on the test: four leaves, 200 slow trees. The search worked, for once. But rerun with four
different random seeds, the defaults' validation score ranges from 0.00\% to 0.17\%, and a search that compared the
candidates against one draw of the defaults could have called anything between the two a success. \omterm{def:ml:trees-and-boosting:boosting}{Gradient boosting} is
the workhorse of the field on tabular data; this chapter builds it from trees, measures how it overfits on market data,
and tunes it with the noise in view.

\section{Trees}

\begin{definition}[Decision tree]\label{def:ml:trees-and-boosting:tree}
A \emph{decision tree}\index{decision tree} partitions the feature space by a sequence of binary splits $x_j\le c$,
chosen greedily to reduce the loss most (for regression, the sum of squared errors), and predicts in each final cell,
or leaf, the mean target of its training rows. Its capacity is set by its depth and by the minimum number of rows per
leaf.
\end{definition}

\begin{omfigure}
\begin{tikzpicture}[font=\small,level distance=1.25cm,level 1/.style={sibling distance=5.2cm},level 2/.style={sibling distance=2.6cm},
  every node/.style={align=center},split/.style={draw=omDef,fill=omDef!10,rounded corners=2pt,inner sep=3pt},
  leaf/.style={draw=omThm,fill=omThm!10,rounded corners=2pt,inner sep=3pt}]
\node[split] {$c_0 \le 0.1$?}
  child {node[split] {$c_2 \le -0.4$?}
    child {node[leaf] {$+0.21\%$\\ \footnotesize 9\,300 rows}}
    child {node[leaf] {$-0.18\%$\\ \footnotesize 17\,100 rows}}}
  child {node[split] {$c_3 \le 0.5$?}
    child {node[leaf] {$+0.12\%$\\ \footnotesize 14\,800 rows}}
    child {node[leaf] {$+0.47\%$\\ \footnotesize 6\,800 rows}}};
\node[font=\footnotesize,anchor=west] at (4.4,-0.1) {yes: left, no: right};
\end{tikzpicture}

\omcaption{A tree of depth two on ranked characteristics, as a schematic: each split sends a stock-month left or right,
each leaf predicts the mean demeaned return of its training rows. The leaf values and counts are illustrative, not
fitted.}\label{fig:ml:trees:tree}
\end{omfigure}

A tree captures interactions without being told (the right branch of \cref{fig:ml:trees:tree} lets characteristic~3
matter only when characteristic~0 is high) and does not care about monotone transformations of the features. Its
weakness is variance: the split points and the leaf means are fitted to noise. On the chapter's panel (300 stocks, 40
ranked characteristics with factor risk, 160 training months, a ceiling of 0.77\% on the 80 test months) a tree of
depth two scores 0.24\% out of sample, depth five $-0.20\%$, depth ten $-3.42\%$.

\section{Bagging and random forests}

\begin{definition}[Bagging, random forest, out-of-bag error]\label{def:ml:trees-and-boosting:bagging}
\emph{Bagging}\index{bagging} (bootstrap aggregation) averages models fitted to bootstrap samples of the \omterm{def:ml:why-financial-machine-learning-is-different:sets}{training set}.
A \emph{random forest}\index{random forest} bags deep trees and, at each split, considers only a random subset of the
features, to decorrelate the trees. The \emph{out-of-bag error}\index{out-of-bag error} scores each training row with
the trees whose bootstrap sample did not contain it.
\end{definition}

\begin{proposition}[What averaging buys]\label{prop:ml:trees:bagging}
If $B$ predictions have the same variance $\sigma^2$ and pairwise correlation $\rho$, their average has variance
$\rho\sigma^2 + (1 - \rho)\sigma^2/B$.
\end{proposition}

\begin{proof}
$\Var(\frac1B\sum_bf_b) = \frac1{B^2}(B\sigma^2 + B(B-1)\rho\sigma^2) = \rho\sigma^2 + (1-\rho)\sigma^2/B$.
\end{proof}

More trees remove only the second term; the first is removed by making the trees less correlated, which is what the
random feature subsets are for. A forest of 40 trees (30\% of the features per split, 300 rows per leaf) scores 0.20\%
on the test months. Its out-of-bag R-squared is 0.42\%, twice as good, and its in-sample one 2.76\%. The out-of-bag
rows are not out of sample on a panel: each shares its month, and so its factor returns, with hundreds of in-bag rows
that have similar characteristics. It is chapter~2's warning about out-of-bag scores, in another guise.

\section{Gradient boosting}

\begin{definition}[Gradient boosting, learning rate]\label{def:ml:trees-and-boosting:boosting}
\emph{Gradient boosting}\index{gradient boosting} builds an additive model $F_M(x) = F_0 + \eta\sum_{m=1}^Mh_m(x)$ one
tree at a time: tree $h_m$ is fitted to the negative gradient of the loss with respect to the current predictions,
$-\partial\ell(y_i, F_{m-1}(x_i))/\partial F$, and added with a factor $\eta\in(0,1]$. The \emph{learning
rate}\index{learning rate} $\eta$ is the step multiplier of an iterative fit: here the shrinkage of each tree, in
gradient descent the size of each step (Book~4, chapter~24). Following the ML convention, $\eta$ denotes it in this
book as a local symbol.
\end{definition}

\begin{proposition}[Boosting as gradient descent in function space]\label{prop:ml:trees:gradient}
For the squared loss $\ell(y, F) = \frac12(y - F)^2$ the negative gradient is the residual $y_i - F_{m-1}(x_i)$, so each
tree fits the current residuals. For any differentiable loss, one step $F_m = F_{m-1} + \eta h_m$ with $h_m$ equal to the
negative gradient at the training points decreases the empirical loss for small enough $\eta$.
\end{proposition}

\begin{proof}
$\partial\frac12(y - F)^2/\partial F = -(y - F)$. For the second claim, $\mathcal L_n(F_{m-1} + \eta h) =
\mathcal L_n(F_{m-1}) - \eta\frac1n\sum_ig_i^2 + O(\eta^2)$ with $g_i$ the gradient at row $i$, when $h = -g$ at the
training points; a tree approximates $-g$ and the same expansion holds with $\sum_ig_ih(x_i) < 0$.
\end{proof}

The Huber loss of Book~4 (chapter~15) is the usual choice when returns have fat tails; LightGBM's histogram-based trees,
grown leaf by leaf, make the method fast enough to fit hundreds of models a day (Ke et al.).

\section{Regularisation, early stopping and monotonic constraints}

\begin{definition}[Early stopping]\label{def:ml:trees-and-boosting:early}
\emph{Early stopping}\index{early stopping} stops an iterative fit (more trees, more epochs) at the iteration where the
loss on a validation block, purged from the training rows, is lowest.
\end{definition}

The number of trees and the \omterm{def:ml:trees-and-boosting:boosting}{learning rate} trade against each other (\cref{fig:ml:trees:early}). At $\eta = 0.1$ the
validation R-squared peaks after 34 trees and collapses afterwards; the stopped model scores 0.05\% on the test months,
the model run to 600 trees $-3.38\%$. At $\eta = 0.01$ the peak comes after 233 trees and is higher and flatter; the
stopped model scores 0.26\% on the test months, and running to 600 trees costs little ($0.03\%$). A small \omterm{def:ml:trees-and-boosting:boosting}{learning rate}
is a regulariser in its own right: each tree moves the fit so little that noise must be fitted slowly.

\begin{omfigure}
\begin{tikzpicture}
\begin{axis}[width=11cm,height=5.6cm,xlabel={number of trees},ylabel={validation $R^2$ (\%)},
  tick label style={font=\small},label style={font=\small},xmin=0,xmax=600,ymin=-1.2,ymax=0.6,grid=major,
  grid style={gray!25},legend columns=2,legend style={font=\footnotesize,at={(0.5,-0.3)},anchor=north,draw=none,
  fill=none,/tikz/every even column/.append style={column sep=8pt}},table/col sep=comma,restrict y to domain=-1.2:1]
\addplot[omThm,very thick] table[x=trees,y=lr01]{figdata/ml/05-trees-and-boosting/early.csv};
\addplot[omDef,very thick] table[x=trees,y=lr001]{figdata/ml/05-trees-and-boosting/early.csv};
\draw[omThm,dashed] (axis cs:34,-1.2) -- (axis cs:34,0.6);
\draw[omDef,dashed] (axis cs:233,-1.2) -- (axis cs:233,0.6);
\legend{learning rate 0.1,learning rate 0.01}
\end{axis}
\end{tikzpicture}

\omcaption{Validation R-squared (40 purged months) of boosted trees as trees are added, at two \omterm{def:ml:trees-and-boosting:boosting}{learning rates}; dashed,
the early-stopping points (34 and 233 trees). The curve at $\eta = 0.1$ leaves the frame after about 340 trees and reaches
$-2.9\%$ at 600. Data: \texttt{ml\_trees.early}.}\label{fig:ml:trees:early}
\end{omfigure}

\begin{definition}[Monotonic constraint]\label{def:ml:trees-and-boosting:monotone}
A \emph{monotonic constraint}\index{monotonic constraint} restricts a model to be non-decreasing (or non-increasing) in a
feature, all others fixed. In a tree ensemble it is enforced by allowing a split on that feature only when the left
child's value does not exceed the right child's (or the reverse), propagated down the tree.
\end{definition}

Microstructure often knows a sign the data are too noisy to show. On firm.tape sessions a larger order-flow imbalance
should never forecast a lower mid price (Book~7, chapter~8). Boosted trees fitted to five-second mid changes on eight
ten-minute sessions (912 non-overlapping labels, ten order-book features) and tested on eight more score an R-squared of
0.333 unconstrained, 0.364 with the seven imbalance and flow features constrained to act positively. Unconstrained, a
push of half a standard deviation in the five-second imbalance lowers the forecast for 7.0\% of the test rows; with the
constraints, for none (\cref{fig:ml:trees:pdp}). The constraint is prior knowledge used as regularisation, and it is also
a guarantee a risk committee can read. (The R-squared are high because the simulated mid trends over a few seconds, a
property of \texttt{firm.tape} Book~7 documents; the comparison is the point, not the level.)

\begin{omfigure}
\begin{tikzpicture}
\begin{axis}[width=10.5cm,height=5.4cm,xlabel={five-second order-flow imbalance (standardised units)},
  ylabel={average forecast (ticks)},tick label style={font=\small},label style={font=\small},grid=major,
  grid style={gray!25},legend columns=2,legend style={font=\footnotesize,at={(0.5,-0.3)},anchor=north,draw=none,
  fill=none,/tikz/every even column/.append style={column sep=8pt}},table/col sep=comma]
\addplot[omThm,very thick,mark=*,mark size=1.2pt] table[x=x,y=free]{figdata/ml/05-trees-and-boosting/pdp.csv};
\addplot[omDef,very thick,mark=square*,mark size=1.2pt] table[x=x,y=mono]{figdata/ml/05-trees-and-boosting/pdp.csv};
\legend{unconstrained,monotone}
\end{axis}
\end{tikzpicture}

\omcaption{Partial dependence of the five-second forecast on the order-flow imbalance (the feature swept over its 2nd to
98th percentile, the others at their observed values), for boosted trees without and with \omterm{def:ml:trees-and-boosting:monotone}{monotonic constraints}. Data:
\texttt{ml\_trees.partial\_dependence} on \texttt{firm.tape}.}\label{fig:ml:trees:pdp}
\end{omfigure}

\section{Tuning under noise}

A random search draws twenty configurations (leaves, \omterm{def:ml:trees-and-boosting:boosting}{learning rate}, trees, rows per leaf, row and column sampling, an
$\ell_2$ penalty) and scores each on the validation block (\cref{fig:ml:trees:tuning}). Their validation R-squared
ranges from $-0.54\%$ to 0.39\%. The winner has four leaves, a \omterm{def:ml:trees-and-boosting:boosting}{learning rate} of 0.02 and 200 trees; retrained with four
seeds it scores 0.36--0.39\% (a standard deviation of 0.014 points), and 0.32\% on the test months, 0.07 points below its
validation score, as chapter~3 predicts for a winner. The defaults, retrained with the same four seeds, score between
0.00\% and 0.17\% (a standard deviation of 0.07 points) and $-0.19\%$ on the test months.

Two lessons come out of it. Complex configurations are noisy as well as bad: the seed-to-seed spread of the defaults is
five times that of the winner, so a single validation run of a complex model is a poor estimate of it, and a search that
rescores its top candidates with several seeds is cheap insurance. And the good region is a plateau: six of the
twenty configurations sit within 0.04 points of the best. Choosing the least complex configuration within one seed
standard deviation of the best (four leaves, 1\,000 rows per leaf, no subsampling) gives 0.35\% on validation and 0.31\%
on test, as good, with a model a seventh as flexible by the measure of \cref{fig:ml:trees:tuning}.

\begin{omfigure}
\begin{tikzpicture}
\begin{axis}[width=10.5cm,height=5.8cm,xlabel={complexity, $\log_{10}(\text{leaves}\times\text{trees}\times\eta/\sqrt{\text{rows per leaf}})$},
  ylabel={validation $R^2$ (\%)},tick label style={font=\small},label style={font=\small},grid=major,grid style={gray!25},
  ymin=-0.65,ymax=0.5,xmin=-0.5,xmax=1.8,xtick={-0.4,0,0.4,0.8,1.2,1.6},legend columns=4,legend style={font=\footnotesize,at={(0.5,-0.3)},anchor=north,
  draw=none,fill=none,/tikz/every even column/.append style={column sep=6pt}},table/col sep=comma]
\addplot[only marks,mark=*,mark size=1.8pt,gray] table[x=complexity,y=score]{figdata/ml/05-trees-and-boosting/tuning.csv};
\addplot[only marks,mark=square*,mark size=2.6pt,omThm] table[x=complexity,y=score]{figdata/ml/05-trees-and-boosting/tuning_best.csv};
\addplot[only marks,mark=diamond*,mark size=3pt,omDef] table[x=complexity,y=score]{figdata/ml/05-trees-and-boosting/tuning_plateau.csv};
\addplot[only marks,mark=triangle*,mark size=2.4pt,omProp] table[x=complexity,y=score]{figdata/ml/05-trees-and-boosting/tuning_default.csv};
\legend{other configurations,best,plateau choice,defaults (4 seeds)}
\end{axis}
\end{tikzpicture}

\omcaption{Twenty random configurations of boosted trees on the validation months, against a crude measure of their
flexibility; the defaults retrained with four seeds. Data: \texttt{ml\_trees.tuning}.}\label{fig:ml:trees:tuning}
\end{omfigure}

\begin{method}[Tuning boosted trees on market data]\label{met:ml:trees:tuning}
\begin{enumerate}
\item Fix a purged validation block (or purged folds) and never tune on the test.
\item Start from a small \omterm{def:ml:trees-and-boosting:boosting}{learning rate} (0.01--0.03) with \omterm{def:ml:trees-and-boosting:early}{early stopping}, few leaves and many rows per leaf.
\item Search randomly over the regularisers; rescore the top few candidates with several seeds.
\item Choose from the plateau (the least flexible within one seed standard deviation of the best), not the peak.
\item Declare \omterm{def:ml:trees-and-boosting:monotone}{monotonic constraints} wherever the economics fixes a sign.
\end{enumerate}
\end{method}

\section{Tutorial: tuning in the fog}\label{tut:ml:trees-and-boosting:tut}

\begin{tutorial}
\textbf{Goal.} Fit trees, a forest and boosted trees to the panel; stop boosting early at two \omterm{def:ml:trees-and-boosting:boosting}{learning rates}; run a
random search and choose from its plateau; constrain order-book features on \texttt{firm.tape}. \textbf{End state:}
the chapter's numbers and \cref{fig:ml:trees:early,fig:ml:trees:pdp,fig:ml:trees:tuning}.
\begin{tutsteps}
\item \textbf{\omterm{def:ml:trees-and-boosting:early}{Early stopping}} through LightGBM's callbacks, keeping the validation curve.
\omcode{code/firm/gbdt/firm_gbdt.py}{38}{48}{Early stopping on a purged validation block.}\label{lst:ml:trees:early}
\item \textbf{The plateau rule and the model dump}: the fitted trees become a dictionary that a NumPy function
evaluates exactly (chapter~26 compiles the same dictionary).
\omcode{code/firm/gbdt/firm_gbdt.py}{63}{98}{Choosing from the plateau; evaluating dumped trees.}\label{lst:ml:trees:dump}
\item \textbf{Run} \texttt{ml\_trees.family()}, \texttt{early()}, \texttt{tuning()}, \texttt{monotone()} and
\texttt{fig\_trees.py}.
\end{tutsteps}
\textbf{What to change next.} Constrain the order-flow imbalance with the wrong sign and measure the cost; rescore every
configuration of the search with four seeds and see whether the winner changes.
\end{tutorial}

\section{Build: the firm's boosted trees}\label{bld:ml:trees-and-boosting:gbdt}

\begin{build}
\bfield{Purpose} One way to fit, stop, constrain, ensemble, choose and export gradient-boosted trees.
\bfield{Interface} \texttt{make(params, monotone, names, seed)}, \texttt{DEFAULTS}, \texttt{fit\_early\_stop(params,
X, y, Xv, yv, max\_trees, patience, monotone, names, seed)}, \texttt{SeedEnsemble(params, seeds, monotone, names)},
\texttt{plateau\_select(scores, se, complexity)}, \texttt{to\_dict(model)}, \texttt{predict\_dict(d, X)}.
\bfield{Rules} Deterministic, single-threaded fits; the validation block is purged by the caller; constraints are
declared by feature name, never by column position; exported dictionaries reproduce predictions exactly, missing values
included.
\bfield{Acceptance tests} \texttt{code/firm/gbdt/tests/}: two fits give identical predictions; the dictionary evaluator
matches LightGBM with missing values; \omterm{def:ml:trees-and-boosting:early}{early stopping} keeps the best iteration and the dump holds only it; a monotone
feature never lowers the prediction; seed ensembles average; the plateau rule picks the simplest near-best candidate.
\bfield{Stretch} Quantile and Huber losses through the same interface; a warm-started refit that adds trees for new
months (chapter~12).
\end{build}

\begin{omsources}
\item L.~Breiman, ``Random forests'', \emph{Machine Learning} 45, 2001.
\item J.~H.~Friedman, ``Greedy function approximation: a gradient boosting machine'', \emph{Annals of Statistics} 29(5),
2001.
\item G.~Ke et al., ``LightGBM: a highly efficient gradient boosting decision tree'', \emph{Advances in Neural Information
Processing Systems} 30, 2017.
\item T.~Chen and C.~Guestrin, ``XGBoost: a scalable tree boosting system'', \emph{KDD}, 2016.
\item L.~Grinsztajn, E.~Oyallon and G.~Varoquaux, ``Why do tree-based models still outperform deep learning on typical
tabular data?'', \emph{NeurIPS Datasets and Benchmarks}, 2022.
\end{omsources}

\section{Exercises}

\begin{exercise}[$\star$]\label{exo:ml:trees-and-boosting:1}
Fifty trees each have prediction variance $\sigma^2$ and pairwise correlation 0.3. What is the variance of their
average? And with infinitely many trees?
\end{exercise}

\begin{exercise}[$\star$]\label{exo:ml:trees-and-boosting:2}
A tree splits 1\,000 rows into two leaves of 400 and 600 rows whose mean targets are $-0.1$ and $+0.2$. The overall
mean is $0.08$. By how much does the split reduce the sum of squared errors?
\end{exercise}

\begin{exercise}[$\star$]\label{exo:ml:trees-and-boosting:3}
Boosting with squared loss starts from $F_0 = 0$; a row has target 1.0 and the first tree predicts 0.8 for it. With
$\eta = 0.1$, what is the row's residual after the first step? What would it be after ten steps if every tree, the first
included, predicted its current residual exactly?
\end{exercise}

\begin{exercise}[$\star\star$]\label{exo:ml:trees-and-boosting:4}
Why is the forest's out-of-bag R-squared (0.42\%) twice its test R-squared (0.20\%) on the panel?
\end{exercise}

\begin{exercise}[$\star\star$]\label{exo:ml:trees-and-boosting:5}
The defaults' validation score over four seeds is 0.17, 0.11, 0.09 and 0.00\%. A candidate scores 0.14\% on one seed.
Is it better than the defaults? What would you do before deciding?
\end{exercise}

\begin{exercise}[$\star\star$]\label{exo:ml:trees-and-boosting:6}
\emph{Find the flaw.} ``We used \omterm{def:ml:trees-and-boosting:early}{early stopping} on the test months to pick the number of trees, and report the test
R-squared at that point.''
\end{exercise}

\begin{exercise}[$\star\star\star$]\label{exo:ml:trees-and-boosting:7}
\emph{Coding.} In \texttt{ml\_trees.monotone}, constrain the five-second order-flow imbalance to act \emph{negatively}
(all other constraints as in the chapter). Report the test R-squared and explain the result.
\end{exercise}

\begin{exercise}[$\star\star\star$]\label{exo:ml:trees-and-boosting:8}
Show that with squared loss and a \omterm{def:ml:trees-and-boosting:boosting}{learning rate} $\eta$, if every tree fitted the residuals exactly, a row's residual after
$M$ trees would be $(1 - \eta)^M$ times its target; explain why this makes \omterm{def:ml:trees-and-boosting:early}{early stopping} and a small \omterm{def:ml:trees-and-boosting:boosting}{learning rate}
substitutes.
\end{exercise}

\section{Problem: Tuning in the Fog}

\begin{problem}[{Weekend problem --- twenty configurations and a plateau}]\label{pb:ml:trees-and-boosting:1}
The chapter's panel: 300 stocks, forty characteristics with factor risk, 160 training months (the last 40 of them as a
purged validation block when tuning), 80 test months.

\medskip\noindent\textbf{Part I --- Single trees and forests.}
\begin{enumerate}
\item What do trees of depth 2, 5 and 10 score out of sample, and why the ordering?
\item What does the forest score out of sample, out of bag and in sample?
\item Why is the out-of-bag score optimistic on a panel?
\item What does \cref{prop:ml:trees:bagging} say more trees can and cannot do?
\end{enumerate}

\medskip\noindent\textbf{Part II --- Boosting and \omterm{def:ml:trees-and-boosting:early}{early stopping}.}
\begin{enumerate}[resume]
\item What do the defaults score on the test months?
\item Where does \omterm{def:ml:trees-and-boosting:early}{early stopping} stop at $\eta = 0.1$ and $0.01$, and what do the stopped models score?
\item What would running to 600 trees have cost at each \omterm{def:ml:trees-and-boosting:boosting}{learning rate}?
\item Why is a small \omterm{def:ml:trees-and-boosting:boosting}{learning rate} a regulariser?
\end{enumerate}

\medskip\noindent\textbf{Part III --- The search.}
\begin{enumerate}[resume]
\item What range of validation scores do the twenty configurations cover?
\item What is the winner, and what does it score with four seeds and on the test months?
\item How noisy are the defaults across seeds compared with the winner?
\item Which configuration does the plateau rule choose, and what does it score?
\end{enumerate}

\medskip\noindent\textbf{Part IV --- Constraints and the verdict.}
\begin{enumerate}[resume]
\item What do \omterm{def:ml:trees-and-boosting:monotone}{monotonic constraints} do to the order-book model's R-squared and to its violations?
\item Why can a constraint improve an out-of-sample score?
\item State the \emph{named result}: the seed standard deviation of the defaults' and of the winner's validation
score against the gain from tuning, and the winner's validation-to-test gap.
\item Which configuration would you put in production, and why?
\item How many seeds would you use to rescore the top candidates of a search?
\item What should the research log record about the search?
\item What would change with ten times more stocks?
\item In one sentence: how should boosted trees be tuned on market data?
\end{enumerate}
\end{problem}

\section{Interview questions}

\begin{interviewq}[$\star$ \iqroles{researcher, mle}]\label{iq:ml:trees-and-boosting:1}
Explain the difference between \omterm{def:ml:trees-and-boosting:bagging}{bagging} and boosting.
\end{interviewq}

\begin{interviewq}[$\star$ \iqroles{mle}]\label{iq:ml:trees-and-boosting:2}
Which \omterm{def:ml:why-financial-machine-learning-is-different:sets}{hyperparameters} of a gradient-boosted model would you tune first on noisy data, and in which direction?
\end{interviewq}

\begin{interviewq}[$\star\star$ \iqroles{researcher}]\label{iq:ml:trees-and-boosting:3}
Why might a \omterm{def:ml:trees-and-boosting:bagging}{random forest}'s \omterm{def:ml:trees-and-boosting:bagging}{out-of-bag error} be misleading on financial panel data?
\end{interviewq}

\begin{interviewq}[$\star\star$ \iqroles{researcher, trader}]\label{iq:ml:trees-and-boosting:4}
What is a \omterm{def:ml:trees-and-boosting:monotone}{monotonic constraint}, and give two features of a trading model you would constrain.
\end{interviewq}

\begin{interviewq}[$\star\star$ \iqroles{mle}]\label{iq:ml:trees-and-boosting:5}
Two boosting configurations differ by 0.03 points of validation R-squared. How do you decide between them?
\end{interviewq}

\begin{interviewq}[$\star\star\star$ \iqroles{researcher, mle}]\label{iq:ml:trees-and-boosting:6}
Derive why boosting with squared loss fits residuals, and what changes with a Huber loss.
\end{interviewq}
